\documentclass[11pt,a4paper]{ctexart}
\usepackage{amsmath,amssymb,mathtools,bm}
\usepackage{geometry}
\usepackage{enumitem}
\usepackage{hyperref}
\usepackage{xcolor}
\usepackage{fancyhdr}
\usepackage{titlesec}
\usepackage{booktabs,longtable,array}

\geometry{left=2.05cm,right=2.05cm,top=2.0cm,bottom=2.2cm}
\setmonofont{Consolas}
\hypersetup{
  colorlinks=true,
  linkcolor=blue!45!black,
  urlcolor=blue!45!black,
  citecolor=blue!45!black,
  pdftitle={数值代数高频题库与预测卷},
  pdfauthor={Codex}
}
\setlist{nosep,leftmargin=2em}
\allowdisplaybreaks
\sloppy
\pagestyle{fancy}
\setlength{\headheight}{14pt}
\fancyhf{}
\fancyhead[L]{数值代数期末备考}
\fancyhead[R]{数值代数高频题库与预测卷}
\fancyfoot[C]{\thepage}
\renewcommand{\headrulewidth}{0.4pt}
\newcommand{\code}[1]{\begingroup\ttfamily\small #1\endgroup}
\titlespacing*{\section}{0pt}{1.4em}{0.5em}
\titlespacing*{\subsection}{0pt}{1.0em}{0.35em}
\titlespacing*{\subsubsection}{0pt}{0.8em}{0.25em}

\title{\bfseries 数值代数高频题库与预测卷}
\author{基于历年卷、非上机作业与课程讲义整理}
\date{2026年6月26日}

\begin{document}
\maketitle
\tableofcontents
\newpage

\phantomsection

\section*{一、高频题库}

\addcontentsline{toc}{section}{一、高频题库}

\phantomsection

\subsection*{1. 范数证明}

\addcontentsline{toc}{subsection}{1. 范数证明}

\textbf{题型}：证明 $\|A\|_F$ 是矩阵范数，或证明 $\|A\|_X=\|X^{-1}AX\|_2$ 是矩阵范数。

\textbf{证据}：历年卷 2022-2025；Homework4；总复习。

\textbf{答题骨架}：

\begin{enumerate}

\item 先证非负、齐次；

\item 再证次乘性；

\item $\|A\|_F$ 用 $tr(A^TA)$ 和 $U^TU=I$；

\item $\|A\|_X$ 直接套 2-范数的性质。

\end{enumerate}

\bigskip\hrule\bigskip

\phantomsection

\subsection*{2. LU / PLU / 全主元}

\addcontentsline{toc}{subsection}{2. LU / PLU / 全主元}

\textbf{题型}：给 3 阶矩阵，求 \code{\detokenize{LU}}、\code{\detokenize{PLU}}，并解 $Ax=b$。

\textbf{证据}：历年卷 2021、2025；Homework2。

\textbf{标准例题}：

\[
A=\begin{pmatrix}1&4&7\\2&5&8\\3&6&10\end{pmatrix},\quad b=(1,1,1)^T.
\]

\textbf{答案}：

\[
L=\begin{pmatrix}1&0&0\\2&1&0\\3&2&1\end{pmatrix},\quad
U=\begin{pmatrix}1&4&7\\0&-3&-6\\0&0&1\end{pmatrix},\quad
x=(-1/3,1/3,0)^T.
\]

一种列主元分解：

\[
P=\begin{pmatrix}0&0&1\\1&0&0\\0&1&0\end{pmatrix},\ 
L=\begin{pmatrix}1&0&0\\1/3&1&0\\2/3&1/2&1\end{pmatrix},\ 
U=\begin{pmatrix}3&6&10\\0&2&11/3\\0&0&-1/2\end{pmatrix}.
\]

\textbf{要点}：

\begin{itemize}

\item 普通 LU 可能失败；

\item 选主元是为了控制增长因子和舍入误差；

\item $PAQ=LU$ 只在全主元题里出现。

\end{itemize}

\bigskip\hrule\bigskip

\phantomsection

\subsection*{3. 正定性与平方根法}

\addcontentsline{toc}{subsection}{3. 正定性与平方根法}

\textbf{题型}：证明 SPD，再做 Cholesky / $LDL^T$。

\textbf{证据}：历年卷 2020、2022、2025；Homework3。

\textbf{标准例题}：

\[
A=\begin{pmatrix}
4&-2&4&2\\
-2&10&-2&-7\\
4&-2&8&4\\
2&-7&4&7
\end{pmatrix},\quad
b=(8,2,16,6)^T.
\]

\textbf{答案}：

\[
L=\begin{pmatrix}
2&0&0&0\\
-1&3&0&0\\
2&0&2&0\\
1&-2&1&1
\end{pmatrix},
\quad
LL^T=A,
\quad
x=(1,2,1,2)^T.
\]

\textbf{常考问法}：

\begin{itemize}

\item 说明为什么 \code{\detokenize{A}} 正定；

\item 写出 Cholesky 递推；

\item 说明 $LDL^T$ 为什么比平方根法更稳妥。

\end{itemize}

\bigskip\hrule\bigskip

\phantomsection

\subsection*{4. 最小二乘与正则方程}

\addcontentsline{toc}{subsection}{4. 最小二乘与正则方程}

\textbf{题型}：写正则方程并用平方根法解。

\textbf{证据}：历年卷 2022、2025；Homework7。

\textbf{标准例题}：

\[
A=\begin{pmatrix}0&-1\\3&2\\0&-2\end{pmatrix},\quad
b=\begin{pmatrix}-1/3\\5/3\\-2/3\end{pmatrix}.
\]

\textbf{答案}：

\[
A^TA=\begin{pmatrix}9&6\\6&9\end{pmatrix},\quad
A^Tb=\begin{pmatrix}5\\5\end{pmatrix},\quad
x=(1/3,1/3)^T.
\]

\textbf{证明套路}：

\[
A^T(Ax-b)=0 \iff A^TAx=A^Tb.
\]

\bigskip\hrule\bigskip

\phantomsection

\subsection*{5. Householder 与 Givens}

\addcontentsline{toc}{subsection}{5. Householder 与 Givens}

\textbf{题型}：构造 Householder，或求 Givens 的 $c,s,\beta$。

\textbf{证据}：历年卷 2020、2024、2025；Homework8。

\textbf{标准例题 1}：

\[
x=(3,1,6,4,2,2)^T.
\]

把尾部子向量 $(1,2,2)^T$ 变成 $(3,0,0)^T$，一组可用的 3 维 Householder 为

\[
H=\begin{pmatrix}
1/3&2/3&2/3\\
2/3&1/3&-2/3\\
2/3&-2/3&1/3
\end{pmatrix}.
\]

\textbf{标准例题 2}：

\[
\begin{pmatrix}c&s\\-s&c\end{pmatrix}
\begin{pmatrix}7\\-1\end{pmatrix}
=\begin{pmatrix}\beta\\\beta\end{pmatrix}.
\]

可取

\[
c=3/5,\quad s=-4/5,\quad \beta=5
\]

或等价符号解。

\textbf{要点}：

\begin{itemize}

\item Householder 是“一次消掉一串”；

\item Givens 是“一次消掉一个”；

\item \code{\detokenize{H}} 正交对称，\code{\detokenize{G}} 正交。

\end{itemize}

\bigskip\hrule\bigskip

\phantomsection

\subsection*{6. Jacobi / Gauss-Seidel / SOR}

\addcontentsline{toc}{subsection}{6. Jacobi / Gauss-Seidel / SOR}

\textbf{题型}：写迭代矩阵，判收敛，求最优松弛因子。

\textbf{证据}：历年卷 2020、2022、2024、2025；Homework10。

\textbf{标准模板}：

\[
A=D-L-U.
\]

Jacobi：

\[
B_J=D^{-1}(L+U).
\]

G-S：

\[
B_{GS}=(D-L)^{-1}U.
\]

SOR：

\[
B_\omega=(D-\omega L)^{-1}[(1-\omega)D+\omega U].
\]

\textbf{参数矩阵}：

\[
A=\begin{pmatrix}1&0&a\\0&1&0\\a&0&1\end{pmatrix}
\]

答案：

\begin{itemize}

\item 正定当且仅当 $|a|<1$

\item Jacobi 收敛当且仅当 $|a|<1$

\item G-S 收敛当且仅当 $|a|<1$

\end{itemize}

\textbf{两个 3 阶高频矩阵}：

\[
A_1=\begin{pmatrix}2&1&1\\1&-1&-1\\1&1&-2\end{pmatrix},\quad
A_2=\begin{pmatrix}1&2&-2\\1&1&-1\\2&-2&1\end{pmatrix}
\]

对应 Jacobi 谱半径：

\[
\rho(B_1)=\sqrt5/2>1,\quad \rho(B_2)=0.
\]

所以 \code{\detokenize{A1}} 不收敛，\code{\detokenize{A2}} 收敛。

\bigskip\hrule\bigskip

\phantomsection

\subsection*{7. 共轭梯度法}

\addcontentsline{toc}{subsection}{7. 共轭梯度法}

\textbf{题型}：证明共轭向量线性无关、写 $A^{-1}$ 展开式、手算 CG。

\textbf{证据}：历年卷 2020；Homework11；总复习。

\textbf{必须会写}：

\[
\alpha_k=\frac{r_k^Tr_k}{p_k^TAp_k},\quad
\beta_k=\frac{r_{k+1}^Tr_{k+1}}{r_k^Tr_k},\quad
p_{k+1}=r_{k+1}+\beta_k p_k.
\]

\textbf{证明关键词}：

\begin{itemize}

\item \code{\detokenize{A}} 对称正定；

\item $p_i^T A p_j = 0$；

\item 左乘 $p_j^T A$ 直接锁死系数；

\item $x_k$ 在 Krylov 子空间里。

\end{itemize}

\bigskip\hrule\bigskip

\phantomsection

\subsection*{8. Hessenberg / QR / Schur / SVD}

\addcontentsline{toc}{subsection}{8. Hessenberg / QR / Schur / SVD}

\textbf{题型}：Householder 化上 Hessenberg，证明 QR 迭代保结构，写 Schur 分解，或考 SVD/广义逆。

\textbf{证据}：历年卷 2023；Homework13；参考书 197-199。

\textbf{答题骨架}：

\begin{itemize}

\item 上 Hessenberg：只对第 \code{\detokenize{k}} 列下面部分做 Householder；

\item QR 迭代：$A_{k+1}=Q_k^T A_k Q_k$；

\item Schur：实数域用实 Schur，复数域用上三角 Schur；

\item SVD：$A=U\Sigma V^T$，$A^+=V\Sigma^+U^T$；

\item $AA^+$ 是到列空间的正交投影。

\end{itemize}

\bigskip\hrule\bigskip

\phantomsection

\section*{二、今年最可能的预测卷}

\addcontentsline{toc}{section}{二、今年最可能的预测卷}

\phantomsection

\subsection*{题 1}

\addcontentsline{toc}{subsection}{题 1}

证明 $\|A\|_X=\|X^{-1}AX\|_2$ 是矩阵范数。

\textbf{考点}：范数定义、相似变换、次乘性。

\textbf{关键结论}：直接借 \code{\detokenize{2}}-范数的四条性质。

\phantomsection

\subsection*{题 2}

\addcontentsline{toc}{subsection}{题 2}

\[
A=\begin{pmatrix}1&4&7\\2&5&8\\3&6&10\end{pmatrix},\quad
b=(1,1,1)^T.
\]

求 \code{\detokenize{LU}}、\code{\detokenize{PLU}}，并解 $Ax=b$。

\textbf{关键结论}：$x=(-1/3,1/3,0)^T$。

\phantomsection

\subsection*{题 3}

\addcontentsline{toc}{subsection}{题 3}

\[
A=\begin{pmatrix}0&-1\\3&2\\0&-2\end{pmatrix},\quad
b=\begin{pmatrix}-1/3\\5/3\\-2/3\end{pmatrix}.
\]

\begin{enumerate}

\item 写出正则方程；

\item 用平方根法解。

\end{enumerate}

\textbf{关键结论}：$A^TA=[[9,6],[6,9]]$，$x=(1/3,1/3)^T$。

\phantomsection

\subsection*{题 4}

\addcontentsline{toc}{subsection}{题 4}

\[
x=(3,1,6,4,2,2)^T,\qquad
\begin{pmatrix}c&s\\-s&c\end{pmatrix}\begin{pmatrix}7\\-1\end{pmatrix}
=\begin{pmatrix}\beta\\\beta\end{pmatrix}.
\]

分别求 Householder 和 Givens。

\textbf{关键结论}：Householder 子块可取使 $(1,2,2)^T \mapsto (3,0,0)^T$；Givens 可取 $c=3/5, s=-4/5, \beta=5$。

\phantomsection

\subsection*{题 5}

\addcontentsline{toc}{subsection}{题 5}

\[
A=\begin{pmatrix}1&0&a\\0&1&0\\a&0&1\end{pmatrix}.
\]

\begin{enumerate}

\item 何时正定？

\item 何时 Jacobi 收敛？

\item 何时 G-S 收敛？

\end{enumerate}

\textbf{关键结论}：三问都等价于 $|a|<1$。

\phantomsection

\subsection*{题 6}

\addcontentsline{toc}{subsection}{题 6}

给出两个 3 阶矩阵，分别写 Jacobi 迭代矩阵并判断收敛性。

\textbf{关键结论}：

\[
\rho(B_1)=\sqrt5/2>1,\quad \rho(B_2)=0.
\]

\phantomsection

\subsection*{题 7}

\addcontentsline{toc}{subsection}{题 7}

证明 $p_1,...,p_n$ 是 A-共轭向量系，则线性无关，并推出

\[
A^{-1}=\sum_{k=1}^n\frac{p_kp_k^T}{p_k^TAp_k}.
\]

\phantomsection

\subsection*{题 8}

\addcontentsline{toc}{subsection}{题 8}

证明 Householder 变换可将一般矩阵化为上 Hessenberg 矩阵，并说明 QR 迭代为何保持 Hessenberg 结构。

\textbf{关键结论}：每步只作用在活动子块上，结构不被破坏。

\phantomsection

\subsection*{题 9}

\addcontentsline{toc}{subsection}{题 9}

写出 Schur 分解定理，并说明其与 QR 迭代的关系。

\textbf{关键结论}：复 Schur 是 $A=QTQ^*$，实 Schur 是 $A=QTQ^T$。

\phantomsection

\subsection*{题 10}

\addcontentsline{toc}{subsection}{题 10}

写出 SVD 及 Moore-Penrose 广义逆的基本性质，并证明 $AA^+$ 是正交投影。

\textbf{关键结论}：$P=AA^+$ 对称幂等，故是投影。

\bigskip\hrule\bigskip

\phantomsection

\section*{三、书后习题加练（谨慎参考）}

\addcontentsline{toc}{section}{三、书后习题加练（谨慎参考）}

这些题只在和讲义/历年卷同向时提高权重，不单独当作“必考原题”。

\phantomsection

\subsection*{1. 第 74-75 页}

\addcontentsline{toc}{subsection}{1. 第 74-75 页}

重点是舍入误差、增长因子、后向误差界、Cholesky/列主元 Gauss 的误差分析。

\phantomsection

\subsection*{2. 第 136 页}

\addcontentsline{toc}{subsection}{2. 第 136 页}

重点是：

\begin{itemize}

\item Jacobi / G-S / SOR 的收敛条件；

\item $\rho(B)$ 与最优 $\omega$；

\item 松弛迭代和 G-S 的关系。

\end{itemize}

\phantomsection

\subsection*{3. 第 158-159 页}

\addcontentsline{toc}{subsection}{3. 第 158-159 页}

重点是：

\begin{itemize}

\item 最速下降法误差估计；

\item 共轭梯度法中的正交/共轭关系；

\item $A^{-1}$ 的共轭向量展开；

\item 用 CG 解正定方程组。

\end{itemize}

\phantomsection

\subsection*{4. 第 197-199 页}

\addcontentsline{toc}{subsection}{4. 第 197-199 页}

重点是：

\begin{itemize}

\item 幂法；

\item 逆迭代；

\item 位移 QR；

\item Hessenberg 结构；

\item Schur/QL/QR 类算法。

\end{itemize}

\bigskip\hrule\bigskip

\phantomsection

\section*{四、最短背诵版}

\addcontentsline{toc}{section}{四、最短背诵版}

\begin{enumerate}

\item \code{\detokenize{LU/PLU}} 会算，会写 \code{\detokenize{P,L,U}}，会解释选主元。

\item $\|.\|_F$、$\|.\|_2$、$\|.\|_X$ 会证明。

\item \code{\detokenize{Householder/Givens}} 会构造，会说正交。

\item \code{\detokenize{Jacobi/G-S/SOR}} 会写迭代矩阵，会用 $\rho(B)<1$。

\item $Cholesky/LDL^T$ 会做，会判正定。

\item 最小二乘会写正则方程、会写 QR 解法。

\item CG 会写三条递推，会证共轭向量展开。

\item \code{\detokenize{Hessenberg/QR/Schur/SVD}} 至少会讲清核心定理和为什么这么用。

\end{enumerate}

\end{document}
